\chapter{多元系的复相平衡与化学平衡\\热力学第三定律}

\section{多元均匀系}

\begin{itemize}
    \item 化学变量：$k$个组元的物质的量$\{N_i\}_{i=1}^k$
    \item 偏摩尔量：由于$N_i$是广延量，
    \begin{gather*}
        U(T, p, N_1, \cdots, N_k) = \sum_{i=1}^k N_i \left(\frac{\partial U}{\partial N_i}\right)_{T, p, \{N_{j\neq i}\}} = \sum_{i=1}^k N_i u_i
        \\
        \intertext{类似有}
        V = \sum_{i=1}^k N_i v_i, \quad S = \sum_{i=1}^k N_i s_i, \quad G = \sum_{i=1}^k N_i \mu_i
    \end{gather*}
    组元的$i$偏摩尔量不一定等于纯$i$组元的摩尔量，因为不同分子间有相互作用
    \item 热力学基本微分方程
    \begin{gather*}
        \d U = T \d S - p \d V + \sum_{i=1}^k \mu_i \d N_i
    \end{gather*}
    定义巨势
    \begin{gather*}
        \Psi = F - G
        \\
        \d \Psi = -S \d T - V \d p - \sum_{i=1}^k N_i \d \mu_i
    \end{gather*}
    \item Gibbs-Duhem关系
    \begin{gather*}
        S \d T - V \d p + \sum_{i=1}^k N_i \d \mu_i = 0
    \end{gather*}
\end{itemize}

\section{多元系的复相平衡}

复相平衡条件：设有两个相$\alpha$和$\beta$，且不存在化学反应
\begin{gather*}
    T^{\alpha} = T^{\beta}
    \\
    p^{\alpha} = p^{\beta}
    \\
    \mu_i^{\alpha} = \mu_i^{\beta}, \quad i = 1, 2, \cdots, k
\end{gather*}

化学反应可以表示为
\begin{gather*}
    \sum_{i=1}^k \nu_i A_i = 0
\end{gather*}
其中$\nu_i$为化学计量数，$A_i$为第$i$组元，化学平衡条件
\begin{gather*}
    \sum_{i=1}^k \nu_i \mu_i = 0
\end{gather*}

Gibbs相律
\begin{gather*}
    f = k + 2 - \sigma
\end{gather*}
$f$为热力学自由度，即独立强度变量的个数，$k$为独立组元数，$2$来自强度量$T$和$p$，$\sigma$为独立共存相数

\section{混合理想气体}

设有$k$个理想气体组元，物态方程
\begin{gather*}
    pV=(N_1 + N_2 + \cdots + N_k)RT
\end{gather*}
Dalton分压定律
\begin{gather*}
    p = \sum_{i=1}^k p_i
    \\
    p_i = \frac{N_i RT}{V} = \frac{N_i}{N} p
\end{gather*}

利用对第$i$组元的半透壁，一侧为混合理想气体，另一侧为纯$i$组元理想气体，平衡条件
\begin{gather*}
    p_i = p'_i
    \\
    \mu_i = \mu'_i
\end{gather*}
得到混合理想气体的Gibbs函数
\begin{gather*}
    G = \sum_{i=1}^k N_i \mu_i = \sum_{i=1}^k N_i \mu'_i = \sum_{i=1}^k N_i RT (\varphi_i(T) + \ln(\frac{N_i}{N} p))
    \\
    \varphi_i(T) = \frac{h_{i0} - T s_{i0}}{RT} + \frac{1}{RT} \int c_{p_i} \d T - \frac{1}{R} \int c_{p_i} \frac{\d T}{T}
    \\
    \intertext{熵}
    S = -\left(\frac{\partial G}{\partial T}\right)_{p, \{N_i\}} = \sum_{i=1}^k N_i \left(\int \frac{c_{p_i}}{T} \d T - R \ln (\frac{N_i}{N} p) + s_{i0}\right)
    \\
    \intertext{内能}
    U = G + TS - pV = \sum_{i=1}^k N_i \left(\int c_{v_i} \d T + u_{i0}\right)
\end{gather*}

化学反应平衡时，
\begin{gather*}
    \sum_{i=1}^k \nu_i (\varphi_i(T) + \ln p_i) = 0
    \\
    \intertext{定压平衡常数}
    K_p(T)= \e^{-\sum_{i=1}^k \nu_i \varphi_i(T)} = \prod_{i=1}^k (p_i)^{\nu_i}
\end{gather*}
比较$\prod_{i=1}^k (p_i)^{\nu_i}$与$K_p(T)$，可以判断反应进行的方向

\section{热力学第三定律}

热力学第三定律是量子效应的宏观表现，需要用量子统计理论解释，有三种表述
\begin{itemize}
    \item Nernst定理：系统的熵在等温过程中的改变随绝对温度趋于零而趋于零
    \begin{gather*}
        \lim_{T \to 0} (\Delta S)_T = 0
    \end{gather*}
    推导：对等温等压过程，
    \begin{gather*}
        \intertext{定义化学亲和势}
        A = -\Delta G
        \\
        H = G - T \pt{G}{T}
        \\
        \intertext{定义反应热}
        Q = -\Delta H
        \\
        Q = A - T \pt{A}{T}
        \\
        \intertext{可得}
        \lim_{T \to 0} Q = \lim_{T \to 0} A
        \\
        \pt{A}{T} = \frac{A-Q}{T}
        \\
        \lim_{T \to 0} \pt{A}{T} = \lim_{T \to 0} \frac{A-Q}{T} = \lim_{T \to 0} \left(\pt{A}{T} - \pt{Q}{T}\right)
        \\
        \Delta S = \pt{A}{T}
        \\
        \intertext{Nernst假设}
        \lim_{T \to 0} \pt{A}{T} = \lim_{T \to 0} \pt{Q}{T}
        \\
        \intertext{从而}
        \lim_{T \to 0} (\Delta S)_T = \lim_{T \to 0} \pt{A}{T} = 0
    \end{gather*}
    定理不适用于冻结的非平衡态
    \item 绝对熵：对于$y$不变过程，
    \begin{gather*}
        S(T, y) = S(T_0, y) + \int_{T_0}^T \frac{C_y}{T} \d T
        \\
        (\Delta S)_T = S(0, y_2) - S(0, y_1) + \int_{0}^{T} \frac{C_{y_2} - C_{y_1}}{T} \d T
        \\
        \lim_{T \to 0} (\Delta S)_T = S(0, y_2) - S(0, y_1) = 0
        \\
        \intertext{取}
        S(0, y) = 0
        \\
        S(T, y) = \int_0^T \frac{C_y}{T} \d T
    \end{gather*}
    由$C_y$随$T$趋于$0$，所有物质在绝对零度时的熵为零，即
    \begin{gather*}
        \lim_{T \to 0} S = 0
    \end{gather*}
    \item 不可达原理：不可能通过有限次的热力学过程将系统的温度降低到绝对零度
\end{itemize}